\begin{question}
\noindent In the diagram, $O$ is the centre of a circle. $TA$ is a tangent to the circle at point $A$. $B$ and $C$ are points on the circumference, and $ABCT$ are positioned so that $OAT$ is a straight tangent-radius configuration, with $ABCD$ a cyclic quadrilateral (where $D$ lies on the circle such that $OD$ is also a radius).

The angle between the tangent $TA$ and the chord $AB$ is $34^\circ$. The angle $BOD = 148^\circ$, where $O$ is the centre and $D$ is a point on the circumference such that $OBD$ forms a triangle with two radii $OB$ and $OD$.

\begin{center}
\begin{tikzpicture}[scale=2.2]
    % Circle centre O
    \coordinate (O) at (0,0);
    \draw (O) circle (1);
    \fill (O) circle (0.4pt);
    \node at (O) [above] {$O$};

    % Points on circle
    \coordinate (A) at (200:1);
    \coordinate (B) at (100:1);
    \coordinate (D) at (-40:1);
    \coordinate (C) at (-110:1);

    % Radii
    \draw (O) -- (A);
    \draw (O) -- (B);
    \draw (O) -- (D);

    % Chords forming cyclic quadrilateral ABCD
    \draw (A) -- (B);
    \draw (B) -- (C);
    \draw (C) -- (D);
    \draw (D) -- (A);
    \draw (A) -- (C);

    % Tangent line at A
    \coordinate (T1) at ($(A)!1.6!90:(O)$);
    \coordinate (T2) at ($(A)!1.6!-90:(O)$);
    \draw (T1) -- (T2);
    \node at (T2) [below left] {$T$};

    % Right angle mark at A between OA and tangent
    \tkzMarkRightAngle[size=.15](O,A,T2);

    % Labels
    \node at (A) [left] {$A$};
    \node at (B) [above right] {$B$};
    \node at (C) [below] {$C$};
    \node at (D) [right] {$D$};

    % Angle markings
    \pic [draw, "$34^\circ$", angle eccentricity=1.5, angle radius=0.5cm] {angle = T2--A--B};
    \pic [draw, "$148^\circ$", angle eccentricity=1.4, angle radius=0.6cm] {angle = B--O--D};

\end{tikzpicture}
\end{center}

\begin{enumerate}
    \item[(a)] Triangle $OBD$ is isosceles.\\
    Show that angle $OBD = 16^\circ$, and state the reason why triangle $OBD$ is isosceles.

    \vspace{4cm}
    \answermarks{2}

    \item[(b)] Calculate the size of angle $BCD$, giving reasons for your answer.

    \vspace{4.5cm}
    \answerequnit{BCD}{$^\circ$}{3}
\end{enumerate}
\end{question}
